\documentclass[11pt]{article}

\usepackage[margin=1in]{geometry}
\usepackage{amsmath,amssymb,mathtools}
\usepackage{bm}
\usepackage[hidelinks]{hyperref}

\newcommand{\bs}[1]{\boldsymbol{#1}}

\title{Mapped Tensor-Product Quadrature on the Simplex via Stick-Breaking}
\author{Sachin Natesh}
\date{}

\begin{document}

\maketitle

\section{Reference simplex and Dirichlet--Jacobi weight}

Let
\begin{equation}
  \Delta_d
  = \left\{x\in\mathbb{R}^d:\ x_i\ge 0,\ \sum_{i=1}^d x_i \le 1\right\}
\end{equation}
be the reference $d$-simplex. We use the Dirichlet-type Jacobi weight
\begin{equation}
  w_{\kappa}(x)
  := C_{\kappa}\,
  \left(\prod_{i=1}^d x_i^{\kappa_i-\frac12}\right)
  \left(1-\sum_{i=1}^d x_i\right)^{\kappa_{d+1}-\frac12},
\end{equation}
where $\kappa\in\mathbb{R}^{d+1}$ and the parameters are chosen so that the weight is integrable.
The normalization constant $C_{\kappa}$ is chosen so that
\begin{equation}
  \int_{\Delta_d} w_{\kappa}(x)\,dx = 1.
\end{equation}

\section{Stick-breaking map and Jacobian}

Introduce variables
\begin{equation}
  t=(t_1,\dots,t_d)\in[0,1]^d
\end{equation}
and define the stick-breaking map
\begin{align}
  x_1 &= t_1,\\
  x_2 &= (1-t_1)t_2,\\
  x_3 &= (1-t_1)(1-t_2)t_3,\\
      &\ \vdots\\
  x_d &= \left(\prod_{j=1}^{d-1}(1-t_j)\right)t_d.
  \label{eq:stick-breaking}
\end{align}
Then $x\in\Delta_d$, and the remaining barycentric coordinate is
\begin{equation}
  x_{d+1}
  := 1-\sum_{i=1}^d x_i
  = \prod_{j=1}^d (1-t_j).
\end{equation}
\noindent
The Jacobian determinant of the transformation is
\begin{equation}
  \left|\frac{\partial x}{\partial t}\right|
  = \prod_{j=1}^d (1-t_j)^{d-j}.
  \label{eq:stick-jac}
\end{equation}

\section{Factorization of the weighted measure}

Substituting the stick-breaking map into the weight and multiplying by the Jacobian gives a product measure on the cube:
\begin{equation}
  w_{\kappa}(x(t))
  \left|\frac{\partial x}{\partial t}\right|
  = C_{\kappa}
  \prod_{j=1}^d
  t_j^{\kappa_j-\frac12}
  (1-t_j)^{\eta_j-\frac12},
  \label{eq:weight-factor}
\end{equation}
where
\begin{equation}
  \eta_j
  := \kappa_{j+1}+\kappa_{j+2}+\cdots+\kappa_{d+1}
     + \frac{d-j}{2}.
  \label{eq:eta-def}
\end{equation}
\noindent
Equivalently, if the normalization is absorbed into the one-dimensional factors, the simplex integral becomes
\begin{equation}
  \int_{\Delta_d} f(x)w_{\kappa}(x)\,dx
  =
  \int_{[0,1]^d}
  f(x(t))
  \prod_{j=1}^d
  \widetilde{w}_j(t_j)\,dt,
  \label{eq:simplex-to-cube}
\end{equation}
with
\begin{equation}
  \widetilde{w}_j(t)
  \propto
  t^{\kappa_j-\frac12}
  (1-t)^{\eta_j-\frac12}.
\end{equation}

\section{Mapped tensor-product quadrature}

For each coordinate $j\in\{1,\dots,d\}$, let
\begin{equation}
  \left\{\left(t_q^{(j)},\omega_q^{(j)}\right)\right\}_{q=1}^{Q_j}
\end{equation}
be a Gauss--Jacobi quadrature rule on $[0,1]$ for the weight
\begin{equation}
  t^{\kappa_j-\frac12}(1-t)^{\eta_j-\frac12}.
\end{equation}
\noindent
Taking the tensor product of these one-dimensional rules gives cube nodes
\begin{equation}
  t^{(q)}
  = \left(t_{q_1}^{(1)},\dots,t_{q_d}^{(d)}\right)
\end{equation}
and tensor-product weights
\begin{equation}
  \Omega^{(q)}
  = \prod_{j=1}^d \omega_{q_j}^{(j)}.
\end{equation}
\noindent
Each cube node is mapped to the simplex by \eqref{eq:stick-breaking}:
\begin{equation}
  x^{(q)} = x\bigl(t^{(q)}\bigr).
\end{equation}
The resulting mapped tensor-product rule is
\begin{equation}
  \int_{\Delta_d} f(x)w_{\kappa}(x)\,dx
  \approx
  \sum_q \Omega^{(q)} f\bigl(x^{(q)}\bigr),
  \label{eq:mapped-quadrature}
\end{equation}
with any overall normalization handled consistently by the one-dimensional Jacobi weights.
\newline
\newline
\noindent
This construction converts a weighted integral on the simplex into a tensor product of one-dimensional Jacobi-weighted integrals on $[0,1]$, while the stick-breaking map supplies the corresponding simplex nodes.

\end{document}
